
\section{Compulsion or Freedom}

Editorial article in “Folkebladet” for September 1, 1897.

A reply to Prof. F. A. Schmidt’s article in “Lutheraneren,” which attempts to prove that the congregations are much freer in the old synods than in the Free Church.

We shall begin by asking whether the representative system does not necessarily entail compulsion, if it is to have any meaning at all. When the representatives of our states meet in Congress at Washington, is it not self-evident that, both by law and by logic, they meet with power and authority to decide the matters entrusted to them, and that the states are compelled to submit to their decisions? The United States as a whole forms, by means of President, Congress, and Supreme Court, a higher unity, which must compel each individual state to bow under the decision of the common authority. If the people do not like these decisions, then they must wait until the next election, in case representatives may then be chosen who will change what is not satisfactory.

The representative system remains true to itself when it is transferred into the Church. It imposes a higher unity over the congregation, and it imposes compulsion upon all those congregations which are not convinced of the benefit and necessity of the resolutions adopted. It often makes no difference whether a congregation’s delegate votes for or against a resolution at the annual meeting; the congregation may therefore be opposed to the resolution, and its implementation in the congregation is a compulsion, not a freedom. The representative system with its ruling annual meeting can, when it is misused in the service of machine politics, even bring matters to such a pass that the body decides such things as hardly a single congregation in the whole body would vote in favor of, if it had opportunity to do so. But because it has been decided by the representative annual meeting, all the congregations have both the shame and the harm from it—very much against their will.

In the Free Church there is no representative system. Therefore there is also no ruling annual meeting. The congregations are altogether free, for each congregation by itself decides what it will. No annual meeting decides anything whatsoever for the congregations. Since there are no representatives, there is absolutely no power or authority in the meetings that are held.

The meetings of the Free Church consist of volunteers; and when these volunteers agree on a matter, they recommend it to the congregations. Then the congregations do with the matter as they will. What power or authority, then, have the congregations lost? If they are for the matter, they work for it; if they are against the matter, they take no part in it. No one can command them, and no one can forbid them to work as they find fitting. But one thing there is no room for in the Free Church: a majority cannot force upon a minority matters with which it wants nothing to do. Neither can a majority hinder a minority from working for a cause that it holds dear.

Is this not biblical? Or how was it that Paul’s Gentile mission was carried on? Was it a representative annual meeting that sent him out? Was it the “majority” that compelled the congregations to support him? It has always been the Lord’s way to carry on his work by the instruments he himself has chosen, and by the servants whom his Spirit has driven to the labor. Does not each congregation by itself have full right and warrant to help whatever cause it will?

Georg Sverdrup
